\chapter{Desarrollo de la Solución}

\section{Aqui comenzamos }

El problema se formula como un problema de asignación y secuenciación de tareas sujeto a múltiples restricciones operativas. El objetivo general consiste en maximizar el número de tareas programadas dentro de un horizonte diario de 8 horas por operario, garantizando la factibilidad respecto a precedencias, disponibilidad de repuestos y cualificación del personal.

\subsection{Variables}

\textbf{Variables de decisión:}

\begin{itemize}
    \item $y_i$: operario asignado a la tarea $i$, donde $y_i \in E(\mathrm{tipo}(i))$.
    \item $\pi_i \in [0,1]$: prioridad asignada a la tarea $i$ para determinar su orden de intento de programación.
\end{itemize}

\textbf{Variables auxiliares utilizadas en la simulación:}

\begin{itemize}
    \item $tpo[o]$: tiempo total utilizado por el operario $o$.
    \item $\mathrm{comp}[ot]$: conjunto de tareas completadas dentro de la orden de trabajo $ot$.
    \item $\mathrm{occ}[ot]$: intervalos ocupados en la orden de trabajo $ot$.
\end{itemize}

\subsection{Restricciones}

\begin{itemize}
    \item \textbf{Aptitud del operario:}  
    $y_i \in E(\mathrm{tipo}(i))$.
    \item \textbf{Disponibilidad de repuesto:}  
    Solo puede programarse una tarea si el repuesto asociado a su orden de trabajo está disponible.
    \item \textbf{Precedencias internas de la misma OT:}  
    \[
        \forall t' \in (\mathcal{P}(\mathrm{tipo}(i)) \cap \text{OT}(i)): \; t' \in \mathrm{comp}[\mathrm{ot}(i)]
    \]
    \item \textbf{Horizon­te máximo diario:}  
    $f_i \leq H$.
    \item \textbf{No solapamiento dentro de la misma OT:}  
    \[
        \neg (s_i < f' \land f_i > s'), \quad \forall [s', f') \in \mathrm{occ}[\mathrm{ot}(i)]
    \]
\end{itemize}

\subsection{Función Objetivo}

La función objetivo maximiza el número de tareas programadas y penaliza violaciones a las restricciones:

\[
\operatorname{fitness} =
100 N_{\text{exec}}
- 10 P_{\text{prec}}
- 10 P_{\text{rep}}
- 15 P_{\text{apt}}
- 12 P_{\text{ot}}
- 0.1 E_{\text{hor}}
- 0.5 \sigma
- 0.1 C_{\max}
\]

donde los términos representan:  
$N_{\text{exec}}$ tareas ejecutadas,  
$P_{\text{prec}}$ violaciones de precedencia,  
$P_{\text{rep}}$ faltas de repuestos,  
$P_{\text{apt}}$ asignaciones a operarios no aptos,  
$P_{\text{ot}}$ solapamientos en la misma orden,  
$E_{\text{hor}}$ excede del horizonte,  
$\sigma$ desviación estándar (balance de carga),  
$C_{\max}$ makespan.

% --------------------------------------------------------

\section{Diseño del Algoritmo Genético}

El algoritmo genético se diseñó para explorar combinaciones de asignación y secuenciación que respeten la estructura operativa del taller. Cada individuo representa un cronograma potencial que es evaluado mediante una simulación determinista.

\subsection{Representación}

Cada individuo se define como un vector de $2n$ dimensiones:

\[
[\, y_1, \dots, y_n \,] \;\Vert\; [\, \pi_1, \dots, \pi_n \,]
\]

\begin{itemize}
    \item Primeros $n$ genes: operario asignado, restringido al conjunto $E(\mathrm{tipo}(i))$.
    \item Últimos $n$ genes: prioridades en $[0,1]$.
\end{itemize}

\subsection{Selección}

Se empleó selección por \textbf{ranking}, adecuada para entornos donde la distribución de aptitud puede ser muy desigual debido a penalizaciones fuertes.

\subsection{Cruce}

Se utilizó cruce de tipo \textbf{single\_point}, apropiado debido a que el cromosoma posee dos mitades homogéneas: asignación y prioridad.

\subsection{Mutación}

Mutación aleatoria sobre el \textbf{15\%} de los genes:

\begin{itemize}
    \item Para asignaciones: se elige aleatoriamente un nuevo operario dentro de $E(\mathrm{tipo}(i))$.
    \item Para prioridades: se asigna un valor continuo uniforme en $[0,1]$.
\end{itemize}

\subsection{Parámetros}

\begin{itemize}
    \item Generaciones: 250
    \item Población: 100 individuos
    \item Padres por generación: 10
    \item Padres preservados: 5
    \item Semilla: 42 (reproducibilidad)
\end{itemize}

\subsection*{Evaluación del Individuo}

La simulación determina:

\begin{itemize}
    \item tareas ejecutadas,
    \item tareas rechazadas por precedencias, repuestos, aptitud o solapamiento,
    \item balance de carga entre operarios,
    \item makespan del sistema.
\end{itemize}

% --------------------------------------------------------

\section{Implementación del Método SPT}

Texto pendiente.

% --------------------------------------------------------

\section{Integración con el Entorno Simulado}

Texto pendiente.

% --------------------------------------------------------

\section{Arquitectura General del Sistema Desarrollado}

Texto pendiente.

